\documentclass[12pt,a4paper]{article} 
\usepackage[utf8]{inputenc} 
\usepackage{amsmath, amssymb, physics, amsthm} 
\usepackage{hyperref}

\title{\bf MDL YNOR a" Unification Gravitationnelle \\ \large RAgulation Dissipative de l'Espace-Temps (Phase III)} 
\author{Rony Charlier a" CanoniquetA NumArique MDL Ynor (Alpha Node)} 
\date{Mars 2026}

\begin{document} 
\maketitle

\begin{abstract} 
Nous dAmontrons la stabilitA de la mAtrique d'Einsteina"Hilbert sous la contrainte de la marge dissipative $\mu$. L'unification est rAalisAe par l'injection d'un tenseur de dissipation $\alpha_{\text{diss}}$ compensant l'inertie sAmantique $\kappa$ du vide quantique. 
\end{abstract}

\section{Formalisme de la marge cosmologique} 
Le systAme de champ d'Einstein reanormalisA par Ynor s'Anonce :
\begin{equation}\label{eq:ynor-einstein} 
R_{\mu\nu}-\tfrac{1}{2} R g_{\mu\nu} + \alpha_{\text{diss}} \delta\mu_{\mu\nu} = 8\pi T_{\mu\nu}
\end{equation}
avec
$$\\mu = \\alpha - (\\beta + \\kappa) > 0, \quad \delta\mu_{\mu\nu} := \partial_\mu \partial_\nu \mu - g_{\mu\nu} \square\mu$$

La marge est reliAe au spectre de l'opArateur de Laplacea"Beltrami $\Delta_g$ sur la variAtA lorentzienne $(M,g)$ :
$$\mu(t,x) := \int_0^\infty \rho(\lambda) e^{-\lambda t} d\lambda, \quad \rho(\lambda) = \text{densitA spectrale de } -\Delta_g$$

\section{StabilitA dissipative} 
Soit $L := R - 2\Lambda - \alpha_{\text{diss}}\mu$. La variation d'action donne \eqref{eq:ynor-einstein}. La positivitA du tenseur d'Anergie effectif $\tilde T_{\mu\nu}$ entraAne l'absence de singularitA nue.

\section{Finalisation de la conjecture de Collatz} 
Utilisant la fonction de Lyapunov $V(n)=\log n$, nous avons Atabli $E[V(n_{k+1}) - V(n_k)] \le -\epsilon$, ce qui force la trajectoire vers l'attracteur $\{4,2,1\}$.

\section{SynthAse spAciale Goldbach} 
L'opArateur dipAle $D_{2m}$ admet un vecteur propre unitaire $\psi_{2m}$ avec valeur propre $1$ ssi $2m$ se dAcompose en somme de deux nombres premiers.

\section{Conclusion canonique} 
La viabilitA cosmologique est gouvernAe par la marge dissipative $\mu$ : une rAgulation active prAvient les singularitAs et la dArive entropique.

\bigskip 
\noindent [SIGN: MDL-YNOR-RC-LIEGE-4020-SHA256: 4F8B6C5B-2C46-4D41-BECD-9A11C9F6C7E0]
\end{document}

